% ═══════════════════════════════════════════════════════════════
% MINT MUSTERLÖSUNG TEMPLATE v1.0
% Für Physik, Chemie, Informatik
% ═══════════════════════════════════════════════════════════════
% DESIGN-PRINZIPIEN:
% - Identisches Layout wie AB, nur mit roten Lösungen
% - Fachfarbe dynamisch (Physik/Chemie/Informatik)
% - SW-Modus für druckerfreundliche Version
% - Bewertungshinweise auf separater Lehrerseite
% ═══════════════════════════════════════════════════════════════

\documentclass[11pt, a4paper]{article}

% ═══════════════════════════════════════════════════════════════
% SW-MODUS TOGGLE
% ═══════════════════════════════════════════════════════════════
\newif\ifswmodus
\swmodusfalse    % Standard: Farbe (auf true setzen für SW)

\usepackage[utf8]{inputenc}
\usepackage[T1]{fontenc}
\usepackage[ngerman]{babel}

% --- SCHRIFTEN ---
\usepackage{helvet}
\renewcommand{\familydefault}{\sfdefault}

\usepackage{microtype}
\usepackage[margin=1.5cm, top=2cm, bottom=1.5cm]{geometry}
\usepackage{enumitem}
\usepackage{tabularx}
\usepackage{booktabs}
\usepackage{array}
\usepackage{multicol}
\usepackage{tikz}
\usetikzlibrary{calc}

\usepackage{amsmath, amssymb}
\usepackage{siunitx}
\sisetup{locale = DE, per-mode = symbol}

\usepackage{fancyhdr}
\usepackage{lastpage}
\usepackage{needspace}

\usepackage{xcolor}
\usepackage{tcolorbox}
\tcbuselibrary{skins}

% ═══════════════════════════════════════════════════════════════
% FARBEN
% ═══════════════════════════════════════════════════════════════

% Fachfarben (Basisfarben)
\definecolor{physik-color}{HTML}{2c5aa0}
\definecolor{chemie-color}{HTML}{2a7a4b}
\definecolor{informatik-color}{HTML}{b35c00}

% Strukturfarben
\definecolor{hellgrau-color}{HTML}{F5F5F5}
\definecolor{mittelgrau-color}{HTML}{888888}
\definecolor{dunkelgrau-color}{HTML}{444444}
\definecolor{rahmen-color}{HTML}{CCCCCC}
\definecolor{loesung-color}{HTML}{C62828}

% SW-Farben (druckerfreundlich)
\definecolor{sw-dark}{HTML}{000000}
\definecolor{sw-gray}{HTML}{666666}
\definecolor{sw-white}{HTML}{FFFFFF}

% ═══════════════════════════════════════════════════════════════
% VARIABLEN (ANPASSEN!)
% ═══════════════════════════════════════════════════════════════

\newcommand{\fachid}{physik}                % physik | chemie | informatik
\newcommand{\fach}{Physik}
\newcommand{\thema}{Thema}
\newcommand{\klasse}{10}
\newcommand{\stunde}{01}
\newcommand{\maxpunkte}{25}
\newcommand{\datum}{\today}

% ═══════════════════════════════════════════════════════════════
% BEDINGTE FARBZUWEISUNG
% ═══════════════════════════════════════════════════════════════

\ifswmodus
    \colorlet{fachfarbe}{sw-dark}
    \colorlet{hellgrau}{sw-white}
    \colorlet{mittelgrau}{sw-gray}
    \colorlet{dunkelgrau}{sw-dark}
    \colorlet{rahmen}{sw-gray}
    \colorlet{loesung}{sw-dark}
\else
    \colorlet{fachfarbe}{\fachid-color}
    \colorlet{hellgrau}{hellgrau-color}
    \colorlet{mittelgrau}{mittelgrau-color}
    \colorlet{dunkelgrau}{dunkelgrau-color}
    \colorlet{rahmen}{rahmen-color}
    \colorlet{loesung}{loesung-color}
\fi

% ═══════════════════════════════════════════════════════════════
% KOPF- UND FUSSZEILE
% ═══════════════════════════════════════════════════════════════

\pagestyle{fancy}
\fancyhf{}
\renewcommand{\headrulewidth}{0.4pt}
\renewcommand{\footrulewidth}{0pt}
\lhead{\small \fach{} Klasse \klasse{} | \thema{} \textcolor{loesung}{\textbf{-- MUSTERLÖSUNG}}}
\rhead{\small Name: \underline{\hspace{3.5cm}}}
\cfoot{\small\textcolor{mittelgrau}{Seite \thepage{} von \pageref{LastPage}}}

% ═══════════════════════════════════════════════════════════════
% BOXEN
% ═══════════════════════════════════════════════════════════════

% Merkkasten (farbiger Rahmen oben)
\newtcolorbox{merkkasten}[1][]{
    colback=hellgrau,
    colframe=hellgrau,
    boxrule=0pt,
    arc=0pt,
    left=8pt, right=8pt, top=8pt, bottom=6pt,
    borderline north={2pt}{0pt}{fachfarbe},
    before skip=6pt, after skip=6pt,
    #1
}

% Formelkasten
\newtcolorbox{formelkasten}{
    colback=white,
    colframe=fachfarbe,
    boxrule=1pt,
    arc=0pt,
    left=8pt, right=8pt, top=6pt, bottom=6pt,
    before skip=4pt, after skip=4pt
}

% ═══════════════════════════════════════════════════════════════
% BEFEHLE
% ═══════════════════════════════════════════════════════════════

% Lösung (rot, fett)
\newcommand{\lsg}[1]{\textcolor{loesung}{\textbf{#1}}}

% Lösung unterstrichen (für Lückentexte)
\newcommand{\lsgu}[1]{\textcolor{loesung}{\textbf{\underline{#1}}}}

% Lücke (leer, für Layout-Konsistenz)
\newcommand{\luecke}[1]{\underline{\hspace{#1}}}

% Punktebox
\newcommand{\punktebox}[1]{\hfill\small\textcolor{mittelgrau}{[\lsg{#1}/#1]}}

% Aufgabentitel
\newcommand{\aufgabe}[2]{%
    \needspace{4\baselineskip}%
    \vspace{10pt}%
    \noindent\textbf{\textcolor{fachfarbe}{Aufgabe #1}} #2%
    \vspace{4pt}%
    \nopagebreak[4]%
}

% Operator hervorheben
\newcommand{\operator}[1]{\textbf{\textcolor{fachfarbe}{#1}}}

% Teilaufgaben
\newlist{teil}{enumerate}{2}
\setlist[teil,1]{label=\alph*), leftmargin=1.8em, itemsep=3pt, parsep=0pt, topsep=3pt}
\setlist[teil,2]{label=\roman*), leftmargin=1.8em, itemsep=2pt}

% Einheit (aufrecht in Formeln)
\newcommand{\ein}[1]{\,\mathrm{#1}}

\setlength{\parindent}{0pt}
\setlength{\parskip}{4pt}
\setlength{\columnsep}{1cm}

% ═══════════════════════════════════════════════════════════════
% DOKUMENT
% ═══════════════════════════════════════════════════════════════

\begin{document}

% ═══════════════════════════════════════════════════════════════
% HEADER
% ═══════════════════════════════════════════════════════════════

\begin{center}
    {\Large\bfseries\textcolor{fachfarbe}{\thema}}\\[0.2em]
    {\large\textcolor{loesung}{\textbf{--- MUSTERLÖSUNG ---}}}\\[0.1em]
    {\small Arbeitsblatt | \fach{} Klasse \klasse{} | Stunde \stunde}
\end{center}
\vspace{-0.3em}
\hrule
\vspace{0.6em}

% ═══════════════════════════════════════════════════════════════
% MERKKASTEN (mit Lösungen)
% ═══════════════════════════════════════════════════════════════

\begin{merkkasten}
\textbf{Merke:}

\vspace{0.5em}
\lsg{[Hier steht der vollständige Merksatz]}

\vspace{0.5em}
Formel: $\lsg{v} = \lsg{\dfrac{s}{t}}$
\end{merkkasten}

% ═══════════════════════════════════════════════════════════════
% AUFGABEN MIT LÖSUNGEN
% ═══════════════════════════════════════════════════════════════

\aufgabe{1}{} \punktebox{4}

\operator{Nenne} ...

\lsg{Antwort 1: [Vollständige Lösung]}

\lsg{Antwort 2: [Vollständige Lösung]}

% ---

\aufgabe{2}{} \punktebox{4}

\operator{Ergänze} die Lücken.

\begin{teil}
    \item \lsgu{Lösung a}
    \item \lsgu{Lösung b}
\end{teil}

% ---

\aufgabe{3}{} \punktebox{6}

\operator{Berechne} ...

\textit{Gegeben:} $v = \SI{72}{km/h} = \lsg{\SI{20}{m/s}}$, $t = \SI{5}{s}$

\textit{Gesucht:} $s$

\textit{Formel:} $s = v \cdot t$

\textit{Einsetzen:} $s = \lsg{\SI{20}{m/s}} \cdot \SI{5}{s} = \lsg{\SI{100}{m}}$

\textit{Antwort:} \lsg{Das Auto legt 100 m zurück.}

% ---

\aufgabe{4}{} \punktebox{6}

\operator{Erkläre} ...

\lsg{[Vollständige Erklärung mit allen relevanten Punkten]}

\lsg{[Bewertung: Je 2P für Aspekt A, B, C]}

% ---

\aufgabe{5}{} \punktebox{5}

\operator{Analysiere} ... und \operator{begründe} ...

\lsg{[Analyse: ...]}

\lsg{[Begründung: ...]}

% ═══════════════════════════════════════════════════════════════
% PUNKTE
% ═══════════════════════════════════════════════════════════════

\vfill
\hrule
\vspace{0.4em}

\hfill\textbf{Punkte:} \lsg{\maxpunkte} / \maxpunkte

% ═══════════════════════════════════════════════════════════════
% BEWERTUNGSHINWEISE (NEUE SEITE)
% ═══════════════════════════════════════════════════════════════

\newpage

\begin{center}
    {\Large\textbf{\textcolor{loesung}{Bewertungshinweise für Lehrkräfte}}}
\end{center}

\vspace{1em}

\subsection*{Punktevergabe}

\begin{tabular}{|l|c|l|}
\hline
\textbf{Aufgabe} & \textbf{BE} & \textbf{Kriterien} \\
\hline
1 (Nennen) & 4 & Je 2 BE pro richtige Antwort \\
\hline
2 (Ergänzen) & 4 & Je 2 BE pro richtige Lücke \\
\hline
3 (Berechnen) & 6 & 2 BE Umrechnung, 2 BE Formel/Rechnung, 2 BE Ergebnis+Einheit \\
\hline
4 (Erklären) & 6 & Je 2 BE pro Aspekt (A, B, C) \\
\hline
5 (Analysieren) & 5 & 2 BE Analyse, 3 BE Begründung \\
\hline
\end{tabular}

\vspace{1.5em}

\subsection*{Differenzierte Bewertung}

\textbf{Niveau (*) -- Berufsreife:}
\begin{itemize}
    \item Aufgaben 1--2 vollständig bearbeiten (8 BE)
    \item \textbf{Mindestanforderung:} 6 von 8 BE
\end{itemize}

\textbf{Niveau (**) -- Mittlere Reife:}
\begin{itemize}
    \item Aufgaben 1--4 bearbeiten (20 BE)
    \item \textbf{Mindestanforderung:} 14 von 20 BE
\end{itemize}

\textbf{Niveau (***) -- Gymnasium:}
\begin{itemize}
    \item Alle Aufgaben bearbeiten (25 BE)
    \item \textbf{Mindestanforderung:} 18 von 25 BE
\end{itemize}

\vspace{1.5em}

\subsection*{Häufige Fehler}

\begin{enumerate}
    \item \textbf{Einheiten vergessen:} Ergebnis ohne Einheit $\Rightarrow$ halbe Punktzahl
    \item \textbf{km/h nicht umgerechnet:} Typischer Fehler bei Bewegungsaufgaben
    \item \textbf{Formel falsch umgestellt:} z.B. $t = s \cdot v$ statt $t = \frac{s}{v}$
    \item \textbf{Rechenweg fehlt:} Nur Ergebnis ohne Herleitung $\Rightarrow$ max. 1 BE
\end{enumerate}

\vspace{1.5em}

\subsection*{Nachteilsausgleich}

\textbf{LRS:}
\begin{itemize}
    \item Zeitzugabe: +10\%
    \item Rechtschreibfehler nicht werten
    \item Ggf. vergrößerte Kopie
\end{itemize}

\textbf{Förderbedarf Lernen:}
\begin{itemize}
    \item Nur Aufgaben 1--2 (Niveau *)
    \item Formelkarte als Hilfsmittel erlaubt
    \item Reduzierte BE-Anforderung
\end{itemize}

\end{document}
